
\begin{section}{Proof of Reflection (PoR)}

 PoR is a modification (and associated supporting methods) to existing consensus methods that allows a blockchain to take full advantage of another blockchain's network security (even if consensus methods differ), provided that other blockchain has particular capabilities\footnote{It's very common that a blockchain will have the necessary capabilities, but not all blockchains do.}. A blockchain can use PoR with as many other blockchains as it wishes. Given some blockchain that implements this modification, the difficulty of performing a 51\% attack against it becomes approximately as difficult as performing a 51\% attack against the original blockchain plus \emph{all} the other blockchains that, via PoR, reflect it.

 Additional details are in s3.1 of the whitepaper.

 \begin{subsection}{Prerequisites for PoR}
     In principle, the necessary capabilities (and actions) that some chains,
     \(C_A\) and \(C_B\), must have (and do) in order for \(C_A\) to be
     reflected by \(C_B\) are:

     \begin{enumerate}
         \def\labelenumi{\arabic{enumi}.}
         \tightlist
         \item \label{item:headers-a-in-b}
               The headers of \(C_A\) can be (and are) freely recorded -- promptly and unambiguously -- in \(C_B\);
         \item \label{item:headers-b-in-a}
               The headers of \(C_B\) can be (and are) freely recorded -- promptly and unambiguously -- in \(C_A\); and
         \item \label{item:prove-headers}
               \(C_A\) nodes are able to (and do) promptly prove which of its headers have been recorded in \(C_B\), and have full knowledge of that set of headers.
     \end{enumerate}

     There are multiple well-known methods to achieve \autoref{item:headers-a-in-b} and \autoref{item:headers-b-in-a}, including via specially designed smart contracts or via bespoke protocol enhancements.\footnote{I can provide citations if necessary.}

     A well-known method to facilitate \autoref{item:prove-headers} is via merkle proofs (i.e., one or more branches of a merkle tree) that prove the contents of some relevant state of $C_B$ given a merkle root that is included in $C_B$'s headers. This is often referred to as an \emph{SPV proof}. However, this is not the only method to achieve \autoref{item:prove-headers}. A second method is for every block producer in $C_A$ be capable of verifying the relevant state by following $C_B$'s network rules, e.g., by running a local instance of $C_B$'s client-server software (or equivalent). A third method is for every block producer in $C_A$ to be capable of \emph{guessing} (e.g., by trial and error) the relevant state (given their knowledge of possible options) and subsequently to be capable of efficiently verifying whether the guess was correct or not with certainty (e.g., by cryptographically recreating some known value that is generated from the relevant state). Note, this third method would (probably) not be possible unless $C_B$ had implemented specific features to facilitate it.

     Although the methods to facilitate \autoref{item:prove-headers} may incidentally generate proofs of the relevant state, those proofs \emph{never need to be transmitted between nodes} if certain conditions are met\footnote{See droppable witnesses in whitepaper}.

     Incentivization of block produces to perform the necessary prerequisites are implicit in the construction of PoR (performing these actions is inexpensive and provides their blocks with additional weight, positively impacting their return on investment).

 \end{subsection}
 \begin{subsection}{Required Actions to Facilitate PoR}

     The benefits from \emph{Proof of Reflection} begin as soon as $C_A$ integrates this knowledge into its chain-weighting algorithm, by a method suitably similar to those mentioned in the next section.

 \end{subsection}
 \begin{subsection}{Modification to Existing Consensus Method}

     Implementations of consensus methods use a concept known as \emph{chain-weight} to determine the "heaviest" (or "longest") collection of linked blocks, which is then taken to be the canonical history of the relevant blockchain.
     Methods of consensus will either already have some functions\footnote{Or some collection of code that fulfills an equivalent purpose.} implemented in source code to calculate chain-weight, or be easily modifiable to integrate such functions.
     A generalization of the primary function is shown in \autoref{alg:vanilla-bw}.
     One reason that a blockchain may not explicitly define such a function is because that function would always return 1 (or some trivial result) and has thus been omitted.
     \autoref{alg:refl-1-bw} shows the necessary changes to the calculation of each block's weight.
     Note that \autoref{alg:refl-1-bw} omits a definition of \textsc{WeightOf}.
     Depending on the \emph{reflecting} chain's native currency (or \emph{coin}), one of two definitions should be used. If the blockchains use \emph{the same} coin\footnote{Typically this means that they share the same main/root currency.} then \autoref{alg:weightof-1} or \autoref{alg:weightof-ratio} (or an appropriate variant) should be used.
     However, if the blockchains have \emph{different} coins, then \autoref{alg:weightof-dex} should be used.
     Note that, if \autoref{alg:weightof-dex} is used, then some definition of \textsc{ExchangeRateOfCoinToLocalAt} must be provided (that definition will differ based on consensus algorithm and the available data sources).
     The obvious place to source exchange rate data is via an on-chain decentralized exchange operating between the two blockchains.
     Alternatively, values for given time periods could be hardcoded into the blockchain client-server's source code with new values (for new time periods) added regularly.

     Note that this modification does not depend on participating blockchains using any particular class of consensus methods (e.g., proof of work, proof of stake, etc). Thus PoR can be used to \emph{combine} existing consensus methods. Doing so means that reflected chains gain some of the security qualities of the reflecting chain. For example: reflected proof of stake chains gain the thermodynamic security of reflecting proof of work chains; and reflected proof of work chains gain faster confirmations and possibly low block production variance\footnote{See WP section of similar name}.

     Also note that if a proof stake chain is reflected by a proof of work chain, then the PoW chain provides thermodynamic security to the PoS chain. A class of attacks (for which there were previously no known uncriticized solutions) can be made impossible via hosting the proof of stake set-up and staking procedures on the reflecting PoW chain. This is possible to do with low overhead, since the PoS chain is already tracking parts of the PoW chain's state (so tracking certain other parts is reasonably easy and efficient). By moving these procedures to the PoW chain, it means that an attack on the PoS chain \emph{must} also attack the PoW chain (since the attacks we are guarding against rely on manipulating these procedures); thus the PoS chain is at least as difficult to attack as the PoW chain.

     Also note that PoR can be done \emph{in both directions} between two blockchains (i.e., \emph{mutual} reflection). This increases the security qualities of both chains. When PoR is done in one direction, i.e., Chain A is reflected by Chain B, then an attack on Chain A may require an attack on Chain B \emph{first} -- once Chain B is being successfully attacked, then an attacker can attack Chain A. In the one-way PoR configuration, an attacker can "checkpoint" their progress via sequential attacks. However, when PoR is \emph{mutual}, attacking Chain B requires first attacking Chain A, which first requires attacking Chain B, etc. Thus attacks are more difficult against chains that are mutually reflecting (than configurations where the PoR is one-way). Instead of being able to "checkpoint" their progress, an attacker must attack both chains simultaneously.

 \end{subsection}
 \begin{subsection}{Method of Operation}
     \autoref{fig:por-step3} illustrates the relationship between blockchains where one of those chains (Chain A) is using the other (Chain B) for reflection.
     Note that these diagrams often simplify the depicted blockchains, in that they produce blocks at a consistent and equal rate.
     In reality, there is no requirement that this holds, and PoR works between any two blockchains regardless of their block production rates and is tolerant of production rates that are inconsistent and/or high variance.
     A demonstration of this tolerance is shown in \autoref{fig:por-with-eth}, where a fictional blockchain, Chain A, uses Ethereum for PoR.
     Note that this does not require modification to Ethereum (as it existed in April 2021), so it is currently possible for almost any existing blockchain to implement the necessary changes that initiate and maintain the illustrated relationship.
 \end{subsection}
 \begin{subsection}{Novelty and Inventiveness}
     The following (possibly inexhaustive) list of ideas -- that used in (or based on) PoR -- are novel:

     \begin{enumerate}
         \def\labelenumi{\arabic{enumi}.}
         \tightlist
         \item
               That the block weights of two different blockchains (whether they share a root token or not) are convertible via normalization against the blockchains' respective block rewards\footnote{This avoids involving external market forces, such as hardware availability or the price of electricity.}, and, if necessary, using an exchange rate (such as that which is provided via an on-chain decentralized exchange) to convert denominations.
         \item
               That a blockchain can use the network security of a different (possibly pre-existing and otherwise unassociated) blockchain to enhance its own security via \emph{Proof of Reflection}, provided that a method of converting block-weights exists.
         \item
               That otherwise incompatible consensus methods (such as various proof of work and proof of stake schemes) are able to cooperatively secure each other via PoR.
         \item
               That those otherwise incompatible consensus methods, in doing PoR, inherit some traits of the other consensus method(s); i.e., that PoR is capable of working with all prior consensus methods and enhances their security.
         \item
               That proof of stake chains which are reflected by suitable proof of work chains can move their staking and set-up phases to the reflecting proof of work chain, and, in doing so, thwart a class of attacks for which no previous solution was known. This is a unique feature among blockchain designs, since both the reflecting PoW chain and reflected PoS chain are still individual, stand-alone blockchains. Note that: in no sense does the PoW chain host the PoS chain, as this configuration does not preclude the PoW chain, \emph{simultaneously}, using PoR via reflection in the PoS chain (i.e., the vice-versa configuration).
     \end{enumerate}

 \end{subsection}
 \begin{subsection}{Utility}

     The benefits of PoR (depending on which chains do the reflecting) include: quantitative increase to security, increased confirmation rate (and thus a faster network), the gaining of additional (otherwise unobtainable) security qualities, and more efficient cross-chain transactions with the reflecting chain.

 \end{subsection}
\end{section}

\begin{section}{Construction of The Simplex}

 The Simplex is an emergent construct that forms when a set of blockchains mutually reflect all other blockchains in that set. See \autoref{fig:simplexes}. When in this configuration, for consensus attacks to succeed (e.g., a 51\% attack), an attacker must attack \emph{the simplex as a whole}. A blockchain that is part of a simplex is called a simplex-chain.

 To maximize \emph{both} total throughput of all simplex-chains and the number of simplex-chains in total, the constituent blockchains should devote \nicefrac{1}{2} of their overall capacity\footnote{measured in bytes/s} to reflections (via each miner eagerly including of all new-found headers of other simplex-chains in their blocks).

 Simplexes range in size from 2 chains to tens of thousands of chains (and larger in the future when computational resources are cheaper).

 Additional details are in s3.2 and s3.3 of the whitepaper. This configuration natively has $O(c^2)$ scalability, though this can be improved to $O(c^3)$ or better if the simplex-chains host child blockchains (aka, dapp-chains).

 \begin{subsection}{Requirements for Participating Blockchains}

     Simplex-chains must \emph{mutually} reflect all other simplex-chains and use the modified PoR method. Mutual reflection between two blockchains (which forms a 2-chain simplex) is shown in \autoref{fig:por-step5}.

 \end{subsection}
 \begin{subsection}{Modification to Existing PoR Method}

     Each chain should generalize their reflection capabilities as detailed in \autoref{alg:por-reflected-block-weight}. Note that this includes a cap on each simplex-chain's contribution, which is optional but provides some benefits. Also note that there is more than one way to cap a simplex-chain's contribution, though only one is shown in \autoref{alg:por-reflected-block-weight}.

     For maximum security, each simplex-chain should adjust their weighting function to cap the security contribution of other simplex chains.
     This ensures that real-world variance in block production frequency (of other simplex-chains) does not adversely affect the security of the simplex-chain in question.
     If this is not done, then an attack on one simplex-chain may have amplified effects on the other simplex-chains.
     Note: this change is optional, but beneficial.

 \end{subsection}
 \begin{subsection}{Method of Construction}



 \end{subsection}
 \begin{subsection}{Method of Operation}

 \end{subsection}
 \begin{subsection}{Novelty and Inventiveness}

     \begin{enumerate}
         \def\labelenumi{\arabic{enumi}.}
         \tightlist
         \item
               Generalization of PoR to work over $N$ chains.
         \item
               Capping the security contribution of reflecting chains to secure against some classes of attack.
     \end{enumerate}

 \end{subsection}
 \begin{subsection}{Utility}

     Dramatically improved confirmation rate for simplex-chains over normal blockchains.

     Dramatically improved security for simplex-chains compared to normal blockchains.

 \end{subsection}
\end{section}

\begin{section}{Simplex-Tilings}

 A simplex-tiling is a tiling of simplex-tiles. Simplex-tiles are similar to simplexes. All reflections within a simplex are mutual and \emph{internal}, i.e., with simplex-chains in that simplex. A simplex-\emph{tile} is similar but \emph{not} all reflections are internal.

 Simplex-tilings can be expanded (i.e., more tiles added) in an ad-hoc fashion, i.e., based on network demand. Furthermore, a simplex at $<50\%$ capacity is indistinguishable from a simplex-tiling with only 1 or 2 tiles.

 Additional details are in s3.5 of the whitepaper. This configuration has $O(n)$ scalability, which is the best scalability possible.

 \begin{subsection}{Requirements for Participating Simplex-tiles}

     Simplex-tiles should reserve \nicefrac{1}{8} of their capacity for internal reflections, and \nicefrac{3}{8}\footnote{Simplex-tiles have a property called \emph{valence} which determines how many other simplex-tiles they can do external reflections with. A valence of 3 is the minimum required for a scalable tiling, but smaller and larger valences are possible. The ratios of capacity that a simplex-tile should reserve for internal and external reflections is based on its valence.} of their capacity for \emph{external} reflections, i.e., mutual reflections with \emph{other} tiles. See \autoref{fig:simplex-tiling}.

 \end{subsection}
 \begin{subsection}{Modification to Existing Simplex Parameters}

     The simplex-chains in a simplex-tile should reserve $\frac{1}{2(v + 1)}$, where $v$ is its valence, of its capacity for internal reflections, and $\frac{v}{2(v+1)}$ of its capacity for external reflections. The remaining $\frac{1}{2}$ of capacity is reserved for transactions on that simplex-chain.

 \end{subsection}
 \begin{subsection}{Method of Construction}

     When a simplex-tile is reaching maximum capacity, new blockchains should be created as new simplex-chains in adjacent simplex-tiles. Note that, since simplex-tiles are added in an iterative fashion, all but the first simplex-tile\footnote{Which has $v$ child simplex-tiles} has a parent tile and $v-1$ child tiles. When new simplex-chains are being created on adjacent simplex-tiles, they should only be created in the child-tiles.

 \end{subsection}
 \begin{subsection}{Method of Operation}

     Simplex-chains in a given simplex-tile should mutually reflect all simplex-chains in their own simplex-tile, and all simplex-chains in adjacent simplex-tiles.

 \end{subsection}
 \begin{subsection}{Novelty and Inventiveness}

     \begin{enumerate}
         \def\labelenumi{\arabic{enumi}.}
         \tightlist
         \item
            Creating a high-security and unbounded network of blockchains (organized as a simplex-tiling) via PoR.
         \item
            Creating a scalable method of expansion of this network to remove the upper limit on the number of participating blockchains (and thus remove the upper limit on network throughput and security).
     \end{enumerate}

 \end{subsection}
 \begin{subsection}{Utility}

     $O(n)$ scalability

     Improved security

 \end{subsection}
\end{section}
